\documentclass[11pt,a4paper]{article}
\usepackage[T1]{fontenc}
\usepackage[utf8]{inputenc}
\usepackage{lmodern,amsmath,amssymb,amsthm,mathtools}
\usepackage{geometry,microtype,booktabs,array,longtable}
\usepackage{fancyhdr,fancyvrb,fvextra,needspace,enumitem}
\usepackage{xcolor,hyperref,bookmark}
\geometry{left=25mm,right=25mm,top=25mm,bottom=25mm,headheight=14pt,headsep=9mm}
\definecolor{Ink}{HTML}{17293C}
\definecolor{Link}{HTML}{235577}
\hypersetup{colorlinks=true,linkcolor=Ink,citecolor=Link,urlcolor=Link,
pdftitle={Further Global Closures from Twenty-Five Research Rounds},
pdfauthor={Keio and Shaitetsu},bookmarksnumbered=true}
\urlstyle{same}
\pagestyle{fancy}\fancyhf{}
\fancyhead[L]{\small\textsc{Further global closures}}
\fancyhead[R]{\small\textsc{Rounds 1--25}}
\fancyfoot[C]{\thepage}
\setlength{\parskip}{3pt}
\setlength{\parindent}{1.2em}
\setlength{\emergencystretch}{3em}
\setlist{itemsep=4pt,topsep=5pt}
\setcounter{tocdepth}{1}
\numberwithin{equation}{section}
\newtheorem{theorem}{Theorem}[section]
\newtheorem{lemma}[theorem]{Lemma}
\newtheorem{proposition}[theorem]{Proposition}
\newtheorem{corollary}[theorem]{Corollary}
\theoremstyle{remark}
\newtheorem{remark}[theorem]{Remark}
\newcolumntype{L}[1]{>{\raggedright\arraybackslash}p{#1}}
\allowdisplaybreaks[2]
\begin{document}
\thispagestyle{empty}
\begin{center}
{\small\textsc{Research continuation \quad | \quad 20 November 2026}}\\[16pt]
{\LARGE\bfseries Further Global Closures\\[5pt]
from Twenty-Five Research Rounds}\\[13pt]
{\large Graph Spectra, Trigonometric Inequalities,\\[3pt]
and an Exact Sixth-Power Certificate}\\[18pt]
{\large Keio and Shaitetsu}
\end{center}
\vspace{6pt}
\begin{abstract}
We complete four arguments left unfinished in a twenty-five-round
mathematical research archive. First, we prove Cazenave's inequality
$\mu_1^2\ge4\lambda_2+8$ for every finite simple connected graph
of order at least two. Five explicit Rayleigh witnesses, positive
vertex weights, and a uniform double-broom bound resolve the
remaining subcubic case. Second, we prove the sharp strict lower
bound $-3/8$ for every finite alternating sine sum proposed by Alzer
and Kwong. A fixed rational interval and a choice of one geometric
scale control the full interval and all truncation lengths. Third,
we give the complete shape classification and sharp constants for
the real-parameter quotient
$(p-\sin x/\sin(x/p))/x^2$ on $0<x<\pi/2$. Its nontrivial
transition occurs at $p_*=1.5577320374\ldots$, characterized by a
unique scalar equation. A coefficient majorant controls the entire
previously untreated range $p>3$; parameters with poles are identified
explicitly. Fourth, we independently certify the endpoint of exactly
seventy-eight positive sixth powers. An elementary three-generator
description gives all permanent exceptions, and an exact finite
sumset calculation, six interval extensions, and seventy-three
explicit padding representations prove the global stabilization
statement, using the classical theorem $g(6)=73$.

The fourth result completes the archive's missing verification; the
source paper already discusses that endpoint elsewhere despite
retaining a conjecture label. All proofs, certificate data, and the
standalone verifier are included in this single source file.
\end{abstract}
\noindent\textbf{Keywords.} Graph spectra; Laplacian eigenvalues;
trigonometric inequalities; sharp constants;
geometric scales; positive power series; numerical semigroups;
positive sixth powers; exact finite certificates.

\clearpage
\tableofcontents
\clearpage
\section*{Scope and provenance}
\addcontentsline{toc}{section}{Scope and provenance}

This paper continues the supplied \emph{Math Tournament Rounds 01--25}
archive after the integrated edition \emph{Global Proofs, Proof
Candidates, and Counterexamples}. It targets unfinished arguments.
The four principal results have the following relationship to the
archived checkpoints.

\begin{center}\small
\begin{tabular}{L{.12\textwidth}L{.32\textwidth}L{.46\textwidth}}
\toprule
Checkpoint & Previous boundary & Present completion\\
\midrule
R12 Heisei & Nonbipartite subcubic graphs remained in Cazenave's squared spectral inequality
& Every finite simple connected graph of order at least two, using explicit local witnesses.\\[5pt]
R22 Taisho & Central interval and small-scale reduction for the alternating sine sum
& Every $n\ge0$ and every $x\in[0,\pi]$; sharpness and the equality set of the infinite series.\\[5pt]
R25 Taisho & Sine quotient classified only for $1<p\le3$
& Full parameter range, the unique transition, and exact optimal constants.\\[5pt]
R21 Keio & Finite interval strategy; unverified exceptional set and padding data
& Complete semigroup proof, every finite membership check, all padding witnesses, and the global tail.\\
\bottomrule
\end{tabular}
\end{center}

The graph and trigonometric sections contain proofs without a numerical
sampling premise. The sixth-power section contains an exact computational component:
its finite recurrence is proved, its complete replay program and
witness data are printed, and separate lemmas cover the infinite
range. The numerical value of $p_*$ is illustrative; its defining
equation and uniqueness proof determine the constant exactly.

Each principal proof was separately checked during preparation.
These checks concern the arguments supplied here; they do not
constitute external peer review or an exhaustive determination of
publication priority. In particular, the provenance distinction for
the sixth-power result is recorded in its final subsection.

\paragraph{Compilation and replay.}
This TEX file is self-contained. Compile with pdfLaTeX, repeating
until its references stabilize. No external bibliography, image, or
included manuscript is required. The program in Appendix~\ref{app:certificate}
uses only the Python standard library and contains all of its data;
Python execution is needed only to replay the finite certificate.
\clearpage
\section{Cazenave's adjacency--Laplacian inequality}
\label{caz:section}

Cazenave's Conjecture~6(iii)~\cite{Cazenave2026} concerns the
second-largest adjacency eigenvalue and the largest ordinary
Laplacian eigenvalue of an arbitrary connected simple graph. The
Round~12 Heisei manuscript proved several cases, but left a
nonbipartite subcubic case unresolved. We close that case by five
small subgraph certificates, two positive vertex weightings, and an
explicit bound for every double broom. The certificates are rational
and are included in full below; the proof has no graph-order cap and
no computational dependency.

\begin{theorem}[Cazenave's Conjecture~6(iii)]
\label{caz:main}
Let $G$ be a finite simple connected undirected graph on $n\ge2$
vertices. Write
\[
 \lambda_1\ge\lambda_2\ge\cdots\ge\lambda_n
 \quad\text{and}\quad
 \mu_1\ge\mu_2\ge\cdots\ge\mu_n=0
\]
for the eigenvalues, counted with multiplicity, of its adjacency
matrix $A$ and its ordinary Laplacian $L=D-A$, respectively. Then
\begin{equation}
 \mu_1^2\ge4\lambda_2+8.
 \label{caz:target}
\end{equation}
\end{theorem}

\subsection{Variational facts and exact local certificates}

We use three elementary consequences of the variational
characterization of eigenvalues. First, if $H$ is any subgraph of
$G$, pad $L(H)$ with zero rows and columns on the omitted vertices.
Then
\[
 L(G)-L(H)=\sum_{uv\in E(G)\setminus E(H)}
 (e_u-e_v)(e_u-e_v)^{\mathsf T}
\]
is positive semidefinite. Consequently,
\begin{equation}
 \mu_1(G)\ge\mu_1(H).
 \label{caz:monotonicity}
\end{equation}
The subgraph need not be induced. In particular, the star centered
at a vertex of maximum degree $\Delta$ gives
\begin{equation}
 \mu_1(G)\ge\Delta+1,
 \qquad \lambda_2(G)\le\lambda_1(G)\le\Delta.
 \label{caz:degree-bounds}
\end{equation}
Second, adjacency interlacing gives
\begin{equation}
 \lambda_2(G)\le\lambda_1(G-v)
 \label{caz:interlacing}
\end{equation}
for every vertex $v$. Third, if a vector $w$ has strictly positive
coordinates and $Aw\le B w$ coordinatewise, then
\begin{equation}
 \lambda_1(A)\le B.
 \label{caz:positive-weights}
\end{equation}
For example, multiply the coordinatewise inequality by a
nonnegative Perron eigenvector of $A$.

Here are the five small graphs used in the proof. An edge written
$ij$ denotes the unordered pair $\{i,j\}$; each displayed list is
the complete edge set of the corresponding certificate graph:
\begin{align*}
 E(S)&=\{01,02,03,14,15\},
   &&V(S)=\{0,\ldots,5\},\\
 E(D^+)&=\{01,02,03,12,13,24\},
   &&V(D^+)=\{0,\ldots,4\},\\
 E(H)&=\{01,12,20,03,34,41\},
   &&V(H)=\{0,\ldots,4\},\\
 E(F)&=\{01,12,20,03,14,45\},
   &&V(F)=\{0,\ldots,5\},\\
 E(C^+)&=\{01,12,23,30,04\},
   &&V(C^+)=\{0,\ldots,4\}.
\end{align*}
Thus $S$ is the double star with two adjacent degree-three centers;
$D^+$ is a diamond $K_4-e$ with a pendant edge at a degree-two
vertex; $H$ is the house graph; $F$ is a triangle with pendant paths
of lengths one and two attached at different vertices; and $C^+$
is a four-cycle with a pendant edge.

\begin{lemma}[Five rational Rayleigh certificates]
\label{caz:certificates}
The graphs above satisfy
\[
 \mu_1(S),\ \mu_1(D^+),\ \mu_1(H)>\sqrt{20},
 \qquad
 \mu_1(F),\ \mu_1(C^+)>\sqrt{19}.
\]
\end{lemma}
\begin{proof}
For a nonzero vector $z$ indexed by a graph's vertices, its
Laplacian Rayleigh quotient is
\[
 q=\frac{z^{\mathsf T}Lz}{z^{\mathsf T}z}
   =\frac{\sum_{ij\in E}(z_i-z_j)^2}{\sum_i z_i^2}.
\]
In the following table, coordinates are in the vertex order
$0,1,\ldots$. Every value follows by integer arithmetic from the
edge sets just given.
\begin{center}
\small
\renewcommand{\arraystretch}{1.35}
\begin{tabular}{@{}clcll@{}}
\toprule
Graph & $z$ & $t$ & $q$ & $q^2-t$\\
\midrule
$S$ & $(4,-4,-1,-1,1,1)$ & $20$ & $41/9$ & $61/81$\\
$D^+$ & $(6,6,-10,-5,3)$ & $20$ & $923/206$ & $3209/42436$\\
$H$ & $(5,-5,0,-3,3)$ & $20$ & $157/34$ & $1529/1156$\\
$F$ & $(-4,5,0,1,-2,1)$ & $19$ & $205/47$ & $54/2209$\\
$C^+$ & $(-5,3,-2,3,1)$ & $19$ & $107/24$ & $505/576$\\
\bottomrule
\end{tabular}
\end{center}
All quotients and all displayed differences are positive. Since
$\mu_1\ge q$, the claims follow.
\end{proof}

\subsection{A bound for double brooms of every length}

Let $B_d$ be the tree obtained from a path of length $d\ge2$ by
attaching two pendant leaves to each endpoint. The path endpoints
therefore have degree three.

\begin{lemma}[Uniform double-broom bound]
\label{caz:broom}
For every integer $d\ge2$,
\[
 \mu_1(B_d)>2+\sqrt5.
\]
\end{lemma}
\begin{proof}
Put
\[
 \tau=2+\sqrt5,
 \qquad r=\frac{\sqrt5-1}{2}\in(0,1).
\]
These constants satisfy
\begin{equation}
 \tau=2+r+r^{-1},
 \qquad
 \tau=3+\frac{2}{\tau-1}+r.
 \label{caz:broom-identities}
\end{equation}
Because $B_d$ is bipartite, changing the sign on one color class
conjugates its Laplacian to $Q=D+A$. It therefore suffices to bound
the largest eigenvalue of $Q$.

Number the vertices of its central path $0,1,\ldots,d$. Define a
strictly positive vector $x$ by
\[
 x_i=r^i+r^{d-i}\qquad(0\le i\le d),
\]
and give each pendant leaf the value of its adjacent endpoint
divided by $\tau-1$. The first identity in
\eqref{caz:broom-identities} gives $(Qx)_i=\tau x_i$ at every
internal path vertex. The same equality holds at every leaf. At
the endpoint $0$, the second identity gives
\begin{align*}
 (Qx)_0-\tau x_0
 &=x_1-rx_0\\
 &=r+r^{d-1}-r(1+r^d)\\
 &=r^{d-1}(1-r^2)>0.
\end{align*}
The other endpoint has the same positive residual. Since all
coordinates of $x$ are positive,
$x^{\mathsf T}Qx>\tau x^{\mathsf T}x$. The Rayleigh principle proves
the claim for every $d$.
\end{proof}

\subsection{Proof of the global inequality}

\begin{proof}[Proof of Theorem~\ref{caz:main}]
Write $\mu=\mu_1(G)$ and let $\Delta$ be the maximum degree.
We first dispose of the degrees outside the remaining subcubic
case, including all paths and cycles.

\paragraph{Maximum degree at least four.}
By \eqref{caz:degree-bounds},
\[
 \mu^2-4\lambda_2-8
 \ge(\Delta+1)^2-4\Delta-8
 =\Delta^2-2\Delta-7\ge1
 \qquad(\Delta\ge4).
\]

\paragraph{Maximum degree at most two.}
Connectedness makes $G$ a path or a cycle. On the path $P_n$,
\[
 \lambda_2=2\cos\frac{2\pi}{n+1},
 \qquad
 \mu=2+2\cos\frac{\pi}{n}.
\]
For $n=2$, \eqref{caz:target} is the equality $4=4$.
For $n\ge3$, put $a=\pi/(2n)$ and $b=\pi/(2(n+1))$.
The target becomes
$16\cos^4a\ge16\cos^2(2b)$. Both cosine quantities are positive,
so this is equivalent to
\[
 \sin a\le\sqrt2\sin b.
\]
Since $\sin t/t$ decreases on $(0,\pi)$,
\[
 \frac{\sin a}{\sin b}\le\frac ab
 =\frac{n+1}{n}\le\frac43<\sqrt2,
\]
which proves the path case. For an even cycle, $\mu=4$ and
$\lambda_2\le2$. For an odd cycle, put $c=\cos(\pi/n)>0$.
Then
\[
 \mu=2+2c\ge4c,
 \qquad
 4\lambda_2+8=8\cos(2\pi/n)+8=16c^2,
\]
and the result follows again.

\paragraph{Initial subcubic reduction.}
Assume henceforth $\Delta=3$. If $\lambda_2\le2$, the star bound
$\mu\ge4$ proves \eqref{caz:target}. If $\mu\ge\sqrt{20}$,
then $\lambda_2\le3$ proves it. It remains to consider
\begin{equation}
 \Delta=3,
 \qquad \lambda_2>2,
 \qquad \mu<\sqrt{20}.
 \label{caz:residual}
\end{equation}
By Lemma~\ref{caz:certificates} and subgraph monotonicity, none of
$S,D^+,H$ can occur as a subgraph.

We claim that $G$ contains no diamond, even as a non-induced
subgraph. In a diamond, two vertices already have degree three,
so only its two degree-two tips could have neighbors outside the
diamond. If any outside vertex exists, connectedness supplies an
edge from a tip to an outside vertex, producing $D^+$. If there is
no outside vertex, $G$ is the diamond or $K_4$; deleting a vertex
and applying interlacing gives $\lambda_2\le2$, contrary to
\eqref{caz:residual}. This proves the claim.

Consequently, the triangles in $G$ are vertex-disjoint. Indeed,
two triangles sharing just one vertex would force degree at least
four, and two sharing an edge would contain a diamond. Moreover,
every edge joining two degree-three vertices lies in a triangle:
if its endpoints had no common neighbor, their four other
neighbors would be distinct and would form the forbidden double
star $S$ with that edge.

In particular, an edge leaving a triangle has its outside endpoint
of degree at most two. Otherwise that edge would lie in another
triangle, contradicting vertex-disjointness. The outside endpoint
cannot itself belong to a different triangle, because its two
neighbors there would make its degree three.

\paragraph{The band $\sqrt{19}\le\mu<\sqrt{20}$.}
Set
\[
 B=1+\sqrt3,
 \qquad c=\frac B2,
 \qquad a=\frac3B,
 \qquad 1<a<c.
\]
Give every triangle vertex weight $c$, every non-triangle vertex
of degree three weight $a$, and every remaining vertex weight one.
At a triangle vertex the adjacency sum divided by its own weight
is at most
\[
 2+\frac1c=B.
\]
At a non-triangle degree-three vertex, all neighbors are
non-triangle vertices of degree at most two, so this ratio is
$3/a=B$. At any remaining vertex the degree is at most two and
all neighbor weights are at most $c$, so the ratio is at most
$2c=B$. Equation~\eqref{caz:positive-weights} therefore yields
\begin{equation}
 \lambda_1(G)\le1+\sqrt3.
 \label{caz:first-perron}
\end{equation}
If $\mu\ge\sqrt{19}$, then
\[
 \mu^2\ge19>12+4\sqrt3
 =8+4(1+\sqrt3)\ge8+4\lambda_2.
\]
The strict scalar comparison follows from $7>4\sqrt3$, or
$49>48$. This closes the whole band.

\paragraph{Triangle attachments below $\sqrt{19}$.}
It remains to assume $\mu<\sqrt{19}$, so $F$ and $C^+$ are also
forbidden. We show that no triangle contains two degree-three
vertices. Suppose otherwise that a triangle has vertices $u,v,w$,
where $u,v$ have respective neighbors $a,b$ outside the triangle.
The vertices $a,b$ are distinct, since otherwise there is a
diamond, and both have degree at most two by the preceding
structural reduction.

If $a$ is not a leaf, let $c$ be its other neighbor. If $c$ is a
triangle vertex, the resulting edges give a diamond (or violate
the already specified degrees). If $c=b$, the original triangle
and the path $u,a,b,v$ give the house $H$. If $c$ is outside
$\{u,v,w,a,b\}$, the triangle together with the path $u,a,c$
and the edge $v,b$ gives $F$. Each case is impossible. Thus $a$
is a leaf, and by symmetry $b$ is a leaf.

If $w$ also has degree three, repeat the same argument for the
pair $u,w$ to show that its outside neighbor is a leaf. All edges
leaving this triangle therefore end in leaves. Connectedness
makes $G$ exactly a triangle with two or three pendant leaves.
Deleting one of its degree-three triangle vertices leaves a path
of order at most four and an isolated vertex. Its adjacency
spectral radius is strictly less than two, so
\eqref{caz:interlacing} contradicts $\lambda_2>2$.

Thus each triangle contains at most one degree-three vertex.
Since every edge between degree-three vertices lies in a
triangle, the degree-three vertices form an independent set.
Give each such vertex weight $\sqrt{3/2}$, and all other vertices
weight one. The adjacency ratio is $3/\sqrt{3/2}=\sqrt6$ at a
degree-three vertex, and is at most $2\sqrt{3/2}=\sqrt6$ elsewhere.
Hence
\begin{equation}
 \lambda_1(G)\le\sqrt6.
 \label{caz:second-perron}
\end{equation}

\paragraph{Finding a double broom.}
There must be at least two degree-three vertices. If $v$ were
the only one, every component of $G-v$ would have maximum degree
two. A cycle component could not be connected to $v$ without
creating another degree-three vertex. Thus all components would
be paths, and interlacing would give
$\lambda_2(G)\le\lambda_1(G-v)<2$, a contradiction.

Choose two degree-three vertices $u,v$ at minimum mutual distance
$d$, and a shortest path between them. Independence gives $d\ge2$;
minimality ensures that every internal path vertex has degree two.
Each endpoint has two other neighbors, all off the path by
shortestness. If $d\ge3$, an off-path neighbor cannot be shared by
the two endpoints, since it would give a path of length two.
Thus their four off-path neighbors are distinct, producing $B_d$
as a subgraph.

If $d=2$ and the endpoints share an off-path neighbor, this
neighbor and the internal path vertex form a four-cycle with
$u,v$. The third neighbor of $u$ is outside this four-cycle:
it is distinct from both common neighbors, and it cannot be $v$
because degree-three vertices are independent. The four-cycle
and this additional edge give $C^+$, a contradiction. Therefore
the four off-path neighbors are distinct in this case too, and
again $B_d$ is a subgraph. Any additional edges can simply be
omitted when selecting that subgraph.

Lemma~\ref{caz:broom} and monotonicity now give
$\mu>2+\sqrt5$. Together with \eqref{caz:second-perron},
\[
 \mu^2>(2+\sqrt5)^2
 =9+4\sqrt5
 >8+4\sqrt6
 \ge8+4\lambda_2.
\]
For the middle strict inequality, observe that
\[
 \sqrt6-\sqrt5=\frac1{\sqrt6+\sqrt5}<\frac14.
\]
This proves the remaining case and completes the global proof.
\end{proof}

\begin{remark}[Scope and priority]
\label{caz:scope}
The result is the full inequality of Conjecture~6(iii), for all
permitted graph orders and all connected simple graphs. It is
distinct from the linear-gap conjecture refuted in Round~11 and
from Conjecture~6(ii), which concerns the second-largest Laplacian
eigenvalue. The proof above was developed by extending the
Round~12 reductions. The source statement was rechecked during
this continuation; no assertion of publication priority is made.
\end{remark}

\Needspace{10\baselineskip}
\section{A sharp lower bound for an alternating Ramanujan sum}
\label{sec:ramanujan}

Alzer and Kwong \cite[Section 4, conjecture (27)]{AlzerKwong} define
\begin{equation}\label{eq:Vdef}
 V_n(x)=\sum_{k=0}^{n}\frac{(-1)^k}{3^k}\sin^3(3^k x)
 \qquad(n\in\mathbb Z_{\ge0}).
\end{equation}
Their conjecture asks for the strict bound below on the entire closed
interval. The Round 22 Taisho checkpoint proved the central-interval
case and an identity near $\pi/6$, but left an estimate across all
geometric scales unresolved. We supply that estimate.

\begin{theorem}\label{thm:ramanujan}
For every integer $n\ge0$ and every $x\in[0,\pi]$,
\begin{equation}\label{eq:ramanujan-main}
 V_n(x)>-\frac38.
\end{equation}
The constant $-3/8$ is the largest uniform lower bound, since
\begin{equation}\label{eq:sharpness}
 V_n(\pi/6)=-\frac38+\frac{1}{2\cdot3^n}.
\end{equation}
\end{theorem}

\subsection{An estimate on one fixed interval}

\begin{lemma}\label{lem:one-scale}
Put $a=51/100$ and $f(s)=(1-\cos^3s)/s$. Then
\[
 f(s)>\frac{13}{20}>\frac{3\sqrt3}{8}
 \qquad\left(\frac{51}{100}\le s\le\frac{153}{100}\right).
\]
\end{lemma}

\begin{proof}
On $(0,\pi/2)$ the sign of $f'(s)$ is the sign of
\[
 q(s)=3s\cos^2s\sin s-(1-\cos^3s).
\]
Direct differentiation gives
\[
 q'(s)=3s\cos s\,(1-3\sin^2s).
\]
Thus $q$ strictly increases from $q(0)=0$ up to
$s=\arcsin(1/\sqrt3)$ and then strictly decreases to
$q(\pi/2)=-1$. Consequently $f$ first increases and then decreases.
On any closed subinterval of $(0,\pi/2)$ its minimum is therefore
at an endpoint. The interval $[a,3a]$ lies in $(0,\pi/2)$.

The alternating Taylor bound
\[
 \cos s\le 1-\frac{s^2}{2}+\frac{s^4}{24}
 \qquad(0\le s\le 153/100)
\]
gives the rational estimates
\begin{align*}
 \cos a&\le\frac{698215067}{800000000}<\frac{873}{1000},\\
 \cos(3a)&\le\frac{46300427}{800000000}<\frac16.
\end{align*}
Cosine is positive at both endpoints, so cubing preserves these
inequalities. It follows that
\begin{align*}
 f(a)&>\frac{1-(873/1000)^3}{51/100}
      =\frac{334661383}{510000000}>\frac{13}{20},\\
 f(3a)&>\frac{1-(1/6)^3}{153/100}
      =\frac{5375}{8262}>\frac{13}{20}.
\end{align*}
For completeness, the differences from $13/20$ in these two rational
bounds are $3161383/510000000$ and $47/82620$, respectively.
Finally,
\[
 \left(\frac{13}{20}\right)^2-
 \left(\frac{3\sqrt3}{8}\right)^2=\frac1{1600}>0.
\]
\end{proof}

\subsection{The exact deficit and the scale selection}

\begin{proof}[Proof of Theorem~\ref{thm:ramanujan}]
Every frequency $3^k$ is odd, hence
$V_n(\pi-x)=V_n(x)$. On the central interval
$[\pi/6,5\pi/6]$, the first term is at least $1/8$, whereas
\[
 \sum_{k=1}^n\frac{(-1)^k}{3^k}\sin^3(3^kx)
 \ge-\sum_{k=1}^n3^{-k}
 =-\frac12+\frac1{2\cdot3^n}.
\]
This proves the desired strict inequality there. By symmetry it remains
to consider $x\in[0,\pi/6]$. Write $x=\pi/6-t$, where
$0\le t\le\pi/6$.

For $k\ge1$ the cosine of $3^k\pi/6$ vanishes and its sine is
$(-1)^{k+1}$. Therefore
\begin{equation}\label{eq:exact-deficit}
\begin{split}
 V_n(\pi/6-t)+\frac38
  ={}&\sin^3(\pi/6-t)-\frac18+\frac1{2\cdot3^n}\\
     &+\sum_{k=1}^n\frac{1-\cos^3(3^kt)}{3^k}.
\end{split}
\end{equation}
Every summand in the last sum is nonnegative, regardless of the size
of $3^kt$. If $t=0$, the right side is $1/(2\cdot3^n)$.
The case $n=0$ also follows directly from $V_0(x)=\sin^3x\ge0$.
We may now assume $n\ge1$ and $t>0$.

Let $c=3\sqrt3/8$. On $[0,\pi/6]$, the derivative of
$\sin^3 y$ is increasing, because
\[
 \frac{d^2}{dy^2}\sin^3y=3\sin y\,(2-3\sin^2y)\ge0,
\]
and its largest value is $c$. Integrating this derivative over a
nondegenerate interval ending at $\pi/6$ yields the strict inequality
\begin{equation}\label{eq:initial-deficit}
 \sin^3(\pi/6-t)-\frac18>-ct.
\end{equation}

First suppose $0<t<a=51/100$. If $3^nt<a$, then
\[
 \frac{1}{2\cdot3^n}>\frac{t}{2a}
 =\frac{50}{51}t>ct.
\]
Equations~\eqref{eq:exact-deficit} and \eqref{eq:initial-deficit}
already prove positivity. Otherwise choose the least integer $k\ge1$
such that $3^kt\ge a$. It exists with $k\le n$, and minimality gives
$a\le3^kt<3a$. Lemma~\ref{lem:one-scale} now gives
\[
 \frac{1-\cos^3(3^kt)}{3^k}
 =t f(3^kt)>\frac{13}{20}t>ct.
\]
This single term cancels the possible initial deficit; every remaining
term in \eqref{eq:exact-deficit} is nonnegative.

Finally suppose $a\le t\le\pi/6$. Here
$\pi/3<3a\le3t\le\pi/2$, so $0\le\cos(3t)<1/2$ and
\[
 \frac{1-\cos^3(3t)}3\ge\frac7{24}>\frac18.
\]
Since $\sin^3(\pi/6-t)-1/8\ge-1/8$, the right side of
\eqref{eq:exact-deficit} is again positive. This exhausts all cases.

At $t=0$, the same identity proves \eqref{eq:sharpness}.
Letting $n$ tend to infinity shows that the lower bound cannot be
increased.
\end{proof}

\begin{corollary}\label{cor:infinite-sine}
The uniformly convergent series
\[
 V(x)=\sum_{k=0}^{\infty}(-1)^k3^{-k}\sin^3(3^kx)
\]
satisfies $V(x)\ge-3/8$ on $[0,\pi]$. Equality holds precisely
at $x=\pi/6$ and $x=5\pi/6$.
\end{corollary}

\begin{proof}
Uniform convergence follows from the bound $3^{-k}$ on the absolute
value of the $k$th term. On the central interval the argument above
gives $V(x)+3/8\ge\sin^3x-1/8$, which is strictly positive away
from its two endpoints. On $x=\pi/6-t$ with $0<t<a$, the least
$k$ in the preceding proof exists independently of the truncation
length, so the limiting deficit is strictly greater than
$(13/20-c)t>0$. The large-$t$ estimate also stays strictly positive
in the limit. Symmetry covers the other end of the interval, and
\eqref{eq:sharpness} gives equality at the two stated points.
\end{proof}

\Needspace{10\baselineskip}
\section{Sharp constants for the full real-parameter sine quotient}
\label{sec:sinequotient}

Male\v sevi\'c, Mi\'covi\'c, Mihailovi\'c and Lutovac
\cite[Section 5]{Malesevic} ask for a real-parameter extension,
with sharp constants, of their inequality (15). The integer-parameter
case comes from Hung and Li \cite{HungLi}. We study
\begin{equation}\label{eq:Qdef}
 Q_p(x)=\frac{p-\sin x/\sin(x/p)}{x^2},\qquad 0<x<\pi/2,
\end{equation}
where $p\ne0$ and the denominator is nonzero. For $p>1$, the
Round 25 Taisho checkpoint established a shape classification only
up to $p=3$. The missing unbounded parameter range is treated in
Lemma~\ref{lem:large-p}. We also determine the unique transition
parameter and classify the remaining real values of $p$.

Define
\begin{equation}\label{eq:endpoint-Q}
 \begin{aligned}
 A(p)&=\frac{p^2-1}{6p} &&(p\ne0),\\
 B(p)&=\frac4{\pi^2}\left(p-\csc\frac{\pi}{2p}\right)
       &&\left(p>\frac12\right).
 \end{aligned}
\end{equation}
Let $p_0=\sqrt{7/3}$. A second constant $p_*$ is defined exactly
below; its decimal value is given only for orientation.

\begin{theorem}\label{thm:sine-classification}
There is a unique $\tau\in(0,\pi/2)$ satisfying
\begin{equation}\label{eq:tau}
 2+\tau\cot\tau-\pi\frac{\sin\tau}{\tau}=0.
\end{equation}
Set
\[
 p_* =\frac{\pi}{2\tau}
       =1.5577320374007023\ldots .
\]
Then $p_0<p_*<2$, and for every $p>1$ the following statements
describe the whole interval $0<x<\pi/2$.
\begin{enumerate}[label=(\roman*)]
\item If $1<p\le p_0$, then $Q_p$ is strictly increasing and
\[
 A(p)<Q_p(x)<B(p).
\]
Both constants are sharp.
\item If $p_0<p<p_*$, then $Q_p$ strictly decreases to one
interior minimum and then strictly increases. Let $t_p$ be the unique
solution in $(0,\pi/(2p))$ of
\begin{equation}\label{eq:critical-t}
 2\left(p-\frac{\sin(pt)}{\sin t}\right)
 +t\frac{d}{dt}\left(\frac{\sin(pt)}{\sin t}\right)=0.
\end{equation}
The best lower and upper constants are
\[
 M(p)=\frac{p-\sin(pt_p)/\sin t_p}{p^2t_p^2},
 \qquad \max\{A(p),B(p)\},
\]
and the precise sharp inequality is
\[
 M(p)\le Q_p(x)<\max\{A(p),B(p)\}.
\]
Equality in the lower bound holds exactly at $x=pt_p$.
\item If $p\ge p_*$, then $Q_p$ is strictly decreasing and
\[
 B(p)<Q_p(x)<A(p).
\]
Both constants are sharp. In particular, the original orientation
of inequality (15) holds precisely for $p\ge p_*$ among $p>1$.
\end{enumerate}
\end{theorem}

\subsection{Positive coefficient bounds and a cosine kernel}

For $p>0$ write
\[
 T=\frac{\pi}{2p},\quad u=t^2,\quad
 R_p(t)=\frac{\sin(pt)}{\sin t},\quad
 H_p(u)=\frac{p-R_p(\sqrt u)}{u}.
\]
Then $Q_p(x)=p^{-2}H_p(x^2/p^2)$. For $0\le q\le1$ put
\[
 J_q(u)=\frac{R_q(\sqrt u)-q}{u},
\]
using continuous values at $u=0$ and $J_0=J_1=0$.

\begin{lemma}\label{lem:J}
For $0\le q\le1$, all Taylor coefficients of $J_q$ at zero
are nonnegative, and the series converges for $|u|<\pi^2$.
For $0<q<1$ they are strictly positive. Moreover,
\begin{equation}\label{eq:J-deriv}
 0\le J_q'(u)\le\frac{1}{36(1-u/6)^2}\le\frac1{25}
 \qquad(0\le u\le1).
\end{equation}
\end{lemma}

\begin{proof}
Euler's sine product gives, with locally uniform convergence,
\begin{equation}\label{eq:euler-ratio}
 R_q(\sqrt u)
 =q\prod_{n=1}^{\infty}
  \left(1+\frac{(1-q^2)u}{\pi^2n^2-u}\right).
\end{equation}
All coefficients in each factor are nonnegative. If $0<q<1$,
any single factor contributes a strictly positive coefficient at
every positive degree. Normal convergence on circles of radius
less than $\pi^2$ justifies coefficient passage to the product.

Write $r_j=[u^j]R_q(\sqrt u)$. Coefficientwise, the product
in \eqref{eq:euler-ratio} is bounded by
\[
 q\prod_{n=1}^{\infty}(1-w_nu)^{-1},
 \qquad w_n=\frac1{\pi^2n^2},\qquad \sum_{n\ge1}w_n=\frac16.
\]
The coefficient of degree $j$ in this last product is $q h_j(w)$,
where $h_j$ is the complete homogeneous symmetric sum. Its monomials
occur with coefficient at least one in $(\sum_nw_n)^j$, so
\[
 0\le r_j\le qh_j(w)\le q6^{-j}\le6^{-j}.
\]
The argument can first be made for finitely many factors and then
passed to the limit. Termwise differentiation now yields
\[
 J_q'(u)=\sum_{j\ge2}(j-1)r_j u^{j-2}
 \le\sum_{j\ge2}(j-1)6^{-j}u^{j-2}
 =\frac{1}{36(1-u/6)^2}.
\]
At $u\le1$ this is at most $1/25$.
\end{proof}

\begin{lemma}\label{lem:K}
For $a>0$ define
\[
 K_a(u)=\frac{1-\cos(a\sqrt u)}{u}
       =a^2\int_0^1(1-v)\cos(av\sqrt u)\,dv.
\]
Whenever $0<a\sqrt u<\pi/2$, one has
\[
 K_a'(u)<0,\qquad K_a''(u)>0,\qquad
 -K_a'(u)\ge\frac{a^4}{12\pi}.
\]
\end{lemma}

\begin{proof}
The integral identity follows by direct integration. Differentiation
under this integral is legitimate on compact positive $u$-intervals,
and also extends continuously to $u=0$ by the power series. In
particular,
\[
 -K_a'(u)=\frac{a^3}{2\sqrt u}
  \int_0^1v(1-v)\sin(av\sqrt u)\,dv>0.
\]
The concavity bound $\sin z\ge2z/\pi$ for $0\le z\le\pi/2$
gives the stated lower bound, since
$\int_0^1v^2(1-v)\,dv=1/12$.
For $b>0$ and $z=b\sqrt u$,
\[
 \frac{d^2}{du^2}\cos(b\sqrt u)
 =\frac{b}{4u^{3/2}}(\sin z-z\cos z)>0
 \qquad(0<z<\pi/2),
\]
because the expression in parentheses has derivative $z\sin z>0$
and vanishes at zero. Integrating proves strict convexity of $K_a$.
\end{proof}

\subsection{The unbounded parameter range}

The elementary subtraction identity
\begin{equation}\label{eq:recurrence-R}
 R_p(t)=2\cos((p-1)t)+R_{p-2}(t)
\end{equation}
holds for every real $p$ whenever $\sin t\ne0$.

\begin{lemma}\label{lem:large-p}
For every real $p\ge2$, the function $H_p$ is strictly decreasing
on $0<u<T^2$.
\end{lemma}

\begin{proof}
For $2\le p\le3$, equation~\eqref{eq:recurrence-R} gives
\[
 H_p=2K_{p-1}-J_{p-2}.
\]
Lemma~\ref{lem:K} and coefficient positivity in Lemma~\ref{lem:J}
show that its derivative is strictly negative.

For $p\ge3$, choose the integer $m\ge1$ so that
$p=2m+q$ with $-1<q\le1$. Iterating
\eqref{eq:recurrence-R} yields
\[
 R_p(t)=2\sum_{j=0}^{m-1}\cos(a_jt)+R_q(t),
 \qquad a_j=p-1-2j.
\]
The smallest frequency is $q+1>0$, the largest is $p-1$, and
$a_jt<pt<\pi/2$. If $q\ge0$, then
\[
 H_p=2\sum_{j=0}^{m-1}K_{a_j}-J_q,
\]
whose derivative is strictly negative. If $q=-\beta<0$, oddness
of $R_q$ in $q$ instead gives
\[
 H_p=2\sum_{j=0}^{m-1}K_{a_j}+J_\beta,
 \qquad 0<\beta<1.
\]
Here $u<T^2\le\pi^2/36<1$. The largest frequency alone controls
the possibly positive last derivative:
\[
 H_p'(u)\le-\frac{(p-1)^4}{6\pi}+\frac1{25}
 <-\frac23+\frac1{25}<0.
\]
We used $p-1\ge2$ and $\pi<4$ in the strict middle inequality.
This completes the proof for all real $p\ge2$.
\end{proof}

\subsection{The unique transition and the sharp constants}

For $1<p<2$, equation~\eqref{eq:recurrence-R} gives
\begin{equation}\label{eq:small-p-convex}
 H_p=2K_{p-1}+J_{2-p}.
\end{equation}
Thus $H_p$ is strictly convex on $[0,T^2]$. Its endpoint values
and slopes, obtained by Taylor expansion and direct differentiation,
are
\begin{align}
 H_p(0)&=L:=\frac{p(p^2-1)}6,
 &H_p(T^2)&=U:=\frac{p-\csc T}{T^2},\label{eq:H-endpoints}\\
 H_p'(0)&=-\frac{p(p^2-1)(3p^2-7)}{360},
 &H_p'(T^2)&=\frac{T\csc T\cot T-2(p-\csc T)}{2T^4}.
 \label{eq:H-slopes}
\end{align}
Indeed,
\[
 R_p(t)=p-\frac{p(p^2-1)}6t^2
 +\frac{p(p^2-1)(3p^2-7)}{360}t^4+O(t^6),
\]
and $R_p'(T)=-\csc T\cot T$.

\begin{lemma}\label{lem:transition}
The function $g(t)=2+t\cot t-\pi\sin t/t$ has exactly one zero
$\tau$ in $(0,\pi/2)$. It is negative on $(0,\tau)$ and positive
on $(\tau,\pi/2)$. The number $p_*=\pi/(2\tau)$ satisfies
$\sqrt{7/3}<p_*<2$.
\end{lemma}

\begin{proof}
The sine product and its logarithmic derivative give, for $|t|<\pi$,
\[
 g(t)=3-\pi+\sum_{k\ge1}c_kt^{2k},\qquad
 c_k=\frac{(-1)^{k+1}\pi}{(2k+1)!}
       -\frac{2\zeta(2k)}{\pi^{2k}}.
\]
Here $c_1=\pi/6-1/3>0$, and $c_k<0$ for all $k\ge2$.
For even $k$ this is immediate. For odd $k\ge3$ it suffices to
prove $\pi^{2k+1}<2(2k+1)!$, because $\zeta(2k)>1$.
At $k=3$, the inequalities $\pi<4$ and $\pi^2<10$ give
$\pi^7<4000<10080=2\cdot7!$. Incrementing $k$ multiplies the
left side by $\pi^2<10$ and the right side by
$(2k+2)(2k+3)>10$, proving the claim by induction.

It follows that $G(u)=g(\sqrt u)$ is strictly concave on
$[0,\pi^2/4]$. At the endpoints,
\[
 G(0)=3-\pi<0,\qquad g(\pi/2)=0,
 \qquad g'(\pi/2)=\frac4\pi-\frac\pi2<0.
\]
Thus $G$ is positive just to the left of its right endpoint and
has a zero in its interior. Strict concavity and $G(0)<0$ give
uniqueness and the asserted signs. More explicitly, $G'$ strictly
decreases, so $G$ increases to one maximum and then decreases;
there can be only the stated interior crossing before the endpoint
zero.

The second slope in \eqref{eq:H-slopes} has the sign of $g(T)$,
since
\[
 H_p'(T^2)=\frac{g(T)}{2T^4\sin T}.
\]
At $p=p_0=\sqrt{7/3}$, strict convexity and $H_p'(0)=0$
force $H_p'(T^2)>0$. At $p=2$ one has $H_2=2K_1$, with
$H_2'(T^2)<0$. Hence $p_0<p_*<2$.
\end{proof}

\begin{proof}[Proof of Theorem~\ref{thm:sine-classification}]
For $1<p<2$, strict convexity makes $H_p'$ strictly increasing.
Its initial value in \eqref{eq:H-slopes} is nonnegative precisely
when $p\le p_0$. In this range every interior derivative is
positive. If $p_0<p<p_*$, the initial derivative is negative and
the final derivative is positive by Lemma~\ref{lem:transition}.
There is exactly one critical point, and it is a strict minimum.
If $p_*\le p<2$, the final derivative is nonpositive, so every
interior derivative is strictly negative. Lemma~\ref{lem:large-p}
treats $p\ge2$.

The strictly increasing change of variable $u=x^2/p^2$ transfers
these shape assertions to $Q_p$. The endpoint limits are exactly
$A(p)=L/p^2$ and $B(p)=U/p^2$. Strict convexity in the middle
case puts all interior values below the larger endpoint value,
and differentiation gives equation~\eqref{eq:critical-t} for its
unique minimum. The endpoint limits establish sharpness of all
endpoint constants. For $p<p_*$, either the orientation is reversed
or an interior minimum is smaller than both endpoints, so the
original double inequality cannot hold throughout the interval.
\end{proof}

\subsection{All other real parameters and the domain of finite bounds}

\begin{proposition}\label{prop:other-p}
The remaining parameter values in \eqref{eq:Qdef} have the following
complete description.
\begin{enumerate}[label=(\roman*)]
\item If $1/2<p<1$, then $Q_p$ is strictly decreasing and
$B(p)<Q_p(x)<A(p)$, with sharp endpoint constants.
\item $Q_1$ and $Q_{-1}$ vanish identically.
\item $Q_{1/2}$ is strictly decreasing with
$\sup Q_{1/2}=-1/4$ and $\inf Q_{1/2}=-\infty$.
\item If $0<p<1/2$, the interval contains a pole, and $Q_p$ is
unbounded in both directions on its domain.
\item All negative parameters follow from $Q_{-p}(x)=-Q_p(x)$;
in particular the roles of lower and upper bounds are reversed.
The value $p=0$ does not define \eqref{eq:Qdef}.
\end{enumerate}
\end{proposition}

\begin{proof}
For $1/2\le p<1$, one has $0<t<T\le\pi$, so
$H_p=-J_p$ has strictly negative derivative by
Lemma~\ref{lem:J} and its coefficient argument on the full interval
$0<u<\pi^2$. For $p>1/2$ the limits in
\eqref{eq:H-endpoints} are finite. At $p=1/2$, the limit at
$x\uparrow\pi/2$ is $-\infty$, whereas the limit at zero is
$A(1/2)=-1/4$. For $0<p<1/2$, the denominator $\sin(x/p)$
has a simple zero at the interior point $x=\pi p$, and the
numerator $\sin x$ is positive there. The two sides of this pole
give both infinite signs. Oddness in $p$ and direct substitution
at $p=1$ prove the other assertions.
\end{proof}

\Needspace{10\baselineskip}
\section{Exactly seventy-eight positive sixth powers}
\label{war:section}

Benfield and Lippard ask for the integers representable as a sum of exactly
seventy-eight positive sixth powers in \cite[Conjecture~6.4]{BenfieldLippard}.
The Round~21 Keio checkpoint reduced this question to a finite interval
calculation and a padding construction, but did not independently verify the
permanent exceptional set or the padding representations. We give a complete
certificate for these missing inputs and an elementary description of the
exceptional set. The infinite part of the proof uses Pillai's theorem
$g(6)=73$ \cite{Pillai1940}.

For a subset $E$ of the nonnegative integers, write $jE$ for its $j$-fold
sumset, and put $0E=\{0\}$. Define
\begin{equation}\label{war:def}
 D=\{a^6-1:a\in\mathbb Z_{\geq1}\},\qquad
 S=\bigcup_{j\geq0}jD,\qquad
 B^6=\mathbb Z_{\geq0}\setminus S.
\end{equation}
The translation by $j$ records the requirement that all $j$ powers be
positive: for $n\geq j$,
\begin{equation}\label{war:translation}
 n\text{ is a sum of exactly }j\text{ positive sixth powers}
 \quad\Longleftrightarrow\quad n-j\in jD.
\end{equation}

\begin{theorem}\label{war:main}
The numerical semigroup $S$ satisfies
\[
 S=\langle63,728,117648\rangle=78D.
\]
Its complement has cardinality $355825$ and largest element $711649$.
For every integer $j\geq78$, a positive integer $n$ is a sum of exactly $j$
positive sixth powers if and only if $n\geq j$ and $n-j\notin B^6$.
Moreover, $77D\ne S$, so $78$ is the least stabilization index.
\end{theorem}

\subsection{The permanent exceptional set}

The semigroup computation can be made entirely elementary. In particular,
the proof does not assume the correctness of a remotely stored list of
$355825$ exceptional integers.

\begin{lemma}\label{war:semigroup}
Let $S_0=\langle63,728,117648\rangle$. Its least representatives modulo $63$
are the $63$ integers
\begin{equation}\label{war:apery}
 d(y,z)=728y+117648z,
 \qquad 0\leq y\leq8,\quad 0\leq z\leq6.
\end{equation}
Consequently
\begin{equation}\label{war:gaps}
 \max(\mathbb Z_{\geq0}\setminus S_0)=711649,
 \qquad
 |\mathbb Z_{\geq0}\setminus S_0|=355825.
\end{equation}
Furthermore, $S=S_0$.
\end{lemma}

\begin{proof}
The identities
\[
 9\cdot728=104\cdot63,
 \qquad 7\cdot117648=13072\cdot63
\]
show that any nonnegative combination of the three generators can be written
as $63x+d(y,z)$ with $x\geq0$ and with $y,z$ in the ranges
\eqref{war:apery}. The reduction of $y$ and $z$ increases $x$, so it
preserves its nonnegativity.

The $63$ numbers in \eqref{war:apery} have distinct residues modulo $63$.
Indeed, $728\equiv0\pmod7$ and $117648\equiv-1\pmod7$, so a residue fixes
$z$ uniquely in $\{0,\ldots,6\}$. Once $z$ is fixed, the value of $y$ is
fixed modulo $9$ because $728/7=104\equiv5\pmod9$ is invertible modulo
$9$. Thus the numbers in \eqref{war:apery} are exactly the least
representatives of $S_0$ in their respective residue classes. Denote by $d_r$
the representative congruent to $r$, where $0\leq r<63$. We have the exact
membership criterion
\begin{equation}\label{war:membership}
 m\in S_0\quad\Longleftrightarrow\quad m\geq d_{m\bmod63}
 \qquad(m\in\mathbb Z_{\geq0}).
\end{equation}
The largest missing integer is therefore
\[
 \max_r d_r-63
 =8\cdot728+6\cdot117648-63=711649.
\]
In residue class $r$, the number of missing nonnegative integers is
$(d_r-r)/63$. Summing gives
\begin{align*}
 |\mathbb Z_{\geq0}\setminus S_0|
 &=\frac{1}{63}\sum_{y=0}^{8}\sum_{z=0}^{6}
       (728y+117648z)-\frac{1}{63}\sum_{r=0}^{62}r\\
 &=\frac{7\cdot728\cdot36+9\cdot117648\cdot21}{63}-31
 =355825.
\end{align*}

The three generators of $S_0$ belong to $D$, so $S_0\subseteq S$.
For the reverse inclusion, every $a\geq10$ satisfies
$a^6-1\geq999999>711649$, and hence $a^6-1\in S_0$. Among the remaining
values, $a=1$ gives zero, $a=2,3,7$ give the three generators, and
\begin{align*}
 4^6-1&=65\cdot63,&
 5^6-1&=248\cdot63,&
 6^6-1&=64\cdot728+63,\\
 8^6-1&=4161\cdot63,&
 9^6-1&=730\cdot728.
\end{align*}
Thus every generator of $S$ belongs to $S_0$, proving equality.
\end{proof}

\subsection{A finite sumset certificate}

Set $H=3500000$. Every set $jD\cap[0,H]$ can be computed with exact integer
arithmetic. The key point is that the recurrence below is exact at every
stage, since all summands are nonnegative; it is not an approximation to
the infinite sumset.

\begin{proposition}[Finite certificate]\label{war:finite}
For $V_j=jD\cap[0,H]$, the missing integers in $[0,H]$ have the counts and
largest elements shown in Table~\ref{war:finite-table}. In particular,
\begin{equation}\label{war:finite-consequences}
 [746111,H]\subseteq40D,
 \qquad V_{78}=S\cap[0,H],
 \qquad 715744\notin77D.
\end{equation}
\end{proposition}

\begin{table}[htbp]
\centering
\begin{tabular}{r|r|r}
$j$ & Number of missing integers in $[0,H]$ & Largest missing integer\\
\hline
40 & 360244 & 746110\\
77 & 355826 & 715744\\
78 & 355825 & 711649
\end{tabular}
\caption{Exact finite sixth-power difference sumsets, with $H=3500000$.}
\label{war:finite-table}
\end{table}

\begin{proof}[Exact computational proof]
Starting with $V_0=\{0\}$, compute
\begin{equation}\label{war:recurrence}
 V_{j+1}=
 \left(\bigcup_{\substack{c\in D\\c\leq H}}(V_j+c)\right)\cap[0,H].
\end{equation}
If a nonnegative sum is at most $H$, all its summands and every partial sum
are at most $H$. This proves \eqref{war:recurrence} in both directions.
The finite set of shifts is
\[
 D\cap[0,H]=\{a^6-1:1\leq a\leq12\}.
\]
Represent $V_j$ by the nonnegative integer
$v_j=\sum_{m\in V_j}2^m$, and put $M=2^{H+1}-1$. Then
\eqref{war:recurrence} is implemented exactly by
\[
 v_0=1,\qquad
 v_{j+1}=M\mathbin{\&}
       \left(\mathop{\mathrm{OR}}_{c\in D\cap[0,H]}(v_j\ll c)\right),
\]
where $\ll$, $\mathbin{\&}$ and $\mathrm{OR}$ denote integer left shift,
bitwise AND and bitwise OR. The missing-set encoding is $M\mathbin{\oplus}v_j$,
so its population count gives the second column of
Table~\ref{war:finite-table}, and its bit length minus one gives the third.

The standalone Python certificate printed in the appendix performs these
operations for $1\leq j\leq78$. It also constructs independently the exact
membership encoding specified by \eqref{war:membership} and compares it,
bit for bit, with $v_{78}$. Thus the second assertion in
\eqref{war:finite-consequences} is checked by equality of complete finite
sets, in addition to the count and maximum checks. At row $77$ the code
explicitly checks that bit $715744$ is zero. The execution returns all three
rows in the table and passes every assertion. This finite certificate
establishes \eqref{war:finite-consequences}.
\end{proof}

\subsection{Extending the certified interval}

\begin{lemma}[Interval extension]\label{war:interval}
Suppose that $[L,U]\cap\mathbb Z\subseteq jD$. If $a$ is a positive integer
such that
\begin{equation}\label{war:overlap}
 a^6-(a-1)^6\leq U-L+1,
\end{equation}
then
\[
 [L,U+a^6-1]\cap\mathbb Z\subseteq(j+1)D.
\]
\end{lemma}

\begin{proof}
For each $1\leq b\leq a$, append $b^6-1\in D$ to the representations of
all integers in $[L,U]$. This gives the interval
$[L+b^6-1,U+b^6-1]\cap\mathbb Z$ inside $(j+1)D$.
The consecutive shift differences $b^6-(b-1)^6$ increase with $b$ and are
bounded by \eqref{war:overlap}. Hence these integer intervals overlap or
touch. Their union contains every integer from $L$ to $U+a^6-1$.
\end{proof}

Starting with $j=40$, $L=746111$ and $U=H$, apply
Lemma~\ref{war:interval} six times, using respectively
$a=14,18,24,33,48,74$. Table~\ref{war:bridge-table} records the resulting
intervals. The standalone certificate checks every overlap inequality
\eqref{war:overlap} and every displayed endpoint using integer arithmetic.

\begin{table}[htbp]
\centering
\begin{tabular}{r|r|r|r}
New $j$ & $a$ & Lower endpoint $L$ & New upper endpoint $U$\\
\hline
41 & 14 & 746111 & 11029535\\
42 & 18 & 746111 & 45041758\\
43 & 24 & 746111 & 236144733\\
44 & 33 & 746111 & 1527612701\\
45 & 48 & 746111 & 13758203164\\
46 & 74 & 746111 & 177964693339
\end{tabular}
\caption{Six exact extensions of the row-$40$ interval.}
\label{war:bridge-table}
\end{table}

Since $0\in D$, each $jD$ is contained in $(j+1)D$. Consequently the last
interval in Table~\ref{war:bridge-table} is contained in $78D$.
It overlaps $[0,H]$, so Proposition~\ref{war:finite} proves
\begin{equation}\label{war:bridge}
 S\cap[0,177964693339]\subseteq78D.
\end{equation}

\subsection{A padding certificate for every larger integer}

\begin{lemma}[Padding representations]\label{war:padding}
The integer $P=41253168892$ is a sum of exactly $s$ positive sixth powers
for every integer $5\leq s\leq77$.
\end{lemma}

\begin{proof}[Exact certificate]
The standalone certificate contains $73$ explicit rows, one for each
$s=5,\ldots,77$. A row lists pairs $a:c_a$ of positive integers and is
checked by the two identities
\begin{equation}\label{war:padding-check}
 \sum_a c_a=s,\qquad \sum_a c_a a^6=P.
\end{equation}
For example, the first two rows are
\begin{align*}
 P&=10^6+13^6+45^6+47^6+53^6,\\
 P&=11^6+23^6+25^6+32^6+42^6+57^6.
\end{align*}
All $73$ rows pass \eqref{war:padding-check} exactly. The rows were
constructed independently by forming the finite table of sums of four
positive sixth powers with bases at most $58$, then searching descending
bases for the remaining summands. The proof needs only the explicitly
stored witnesses and their integer sums; it does not rely on the
completeness of that search or on an external data file.
\end{proof}

\begin{lemma}[Infinite tail]\label{war:tail}
Every integer $n>P$ is a sum of exactly $78$ positive sixth powers.
\end{lemma}

\begin{proof}
Pillai's theorem $g(6)=73$ states that every positive integer is a sum of
at most $73$ positive sixth powers \cite{Pillai1940}. Apply it to $n-P>0$.
After omitting zero summands if the theorem is formulated with nonnegative
bases, there are $t$ positive summands, where $1\leq t\leq73$.
By Lemma~\ref{war:padding}, $P$ has a representation with exactly $78-t$
positive summands, because $5\leq78-t\leq77$.
Adding the two representations gives exactly $78$ positive sixth powers.
\end{proof}

\begin{proof}[Proof of Theorem~\ref{war:main}]
Lemma~\ref{war:semigroup} proves the asserted description and cardinality
of $B^6$. The containment $78D\subseteq S$ follows from the definition.
To prove the reverse containment, let $m\in S$ and put $n=m+78$.
If $n\leq P$, then
\[
 0\leq m\leq P-78=41253168814<177964693339,
\]
so \eqref{war:bridge} gives $m\in78D$. If $n>P$,
Lemma~\ref{war:tail} expresses $n$ as exactly $78$ positive sixth powers,
and \eqref{war:translation} again gives $m\in78D$. Therefore $78D=S$.

For every $j\geq78$, the inclusion $0\in D$ implies
$S=78D\subseteq jD\subseteq S$, hence $jD=S$. The representation
criterion follows from \eqref{war:translation}; if $n<j$, positivity of
each summand makes representation impossible.

Finally, Proposition~\ref{war:finite} gives $715744\notin77D$.
But $715744>711649$, so Lemma~\ref{war:semigroup} gives $715744\in S$.
Thus $77D\ne S$. Since $jD\subseteq77D$ for every $j\leq77$, no smaller
index can stabilize the sumsets. This proves sharpness.
\end{proof}

\begin{remark}[Scope and provenance]\label{war:provenance}
The argument is global: the finite calculation is connected to the entire
infinite tail by explicit overlapping intervals and Pillai's unconditional
theorem. No conjectural bound for sufficiently large sixth-power
representations is used.

The source's status descriptions are inconsistent. Its Section~6 labels the
endpoint $78$ as Conjecture~6.4 and states Theorem~6.3 for $j\geq79$,
whereas its Table~1 labels $g(1,6)=78$ a theorem and its Section~11 sketches
stabilization at $78$ \cite{BenfieldLippard}. Accordingly, the contribution
here is an independently reproducible completion of the Round~21 archive's
unverified dependencies, together with the elementary semigroup proof above.
It should not be presented as a priority claim for an entirely new Waring
theorem. The certificate is self-contained apart from the cited classical
theorem: its complete program embeds all $73$ padding witnesses and reads
no external files.
\end{remark}

\Needspace{10\baselineskip}
\section{Boundaries of this continuation}
\label{sec:boundaries}

The four completed arguments above do not settle every unfinished
target in the twenty-five-round archive. The following boundaries
record what remains unproved by this continuation; they are not
claims about the current status of the entire published literature.

\begin{itemize}
\item For the Round 5 Heisei target, Zhi-Wei Sun's Conjecture 5.8(i),
the attempted transfer to a known hypergeometric evaluation still
requires a uniform harmonic congruence. Finite prime checks do not
supply that missing identity.
\item For the Gaussian-binomial infinite coefficientwise
log-concavity target, finite iterates do not establish positivity
at every iteration. No invariant covering arbitrary iteration
depth has been proved here.
\item The q-Fibonomial unimodality target remains outside the
restricted parameter families already handled in the archive.
\item For the Round 22 Reiwa root-monotonicity target, limits and
eventual monotonicity do not establish the statement at every
genus. The initial unrestricted range is not closed here.
\item For the Round 24 Reiwa multiset Schett target, the archived
real-rootedness argument treats repetitions of the smallest label
with the other labels occurring once. The required extension to
other repeated blocks remains unproved here.
\end{itemize}

These items are excluded from the completed-result count. The
present additions consist of three full extensions of previously
partial arguments and one independent completion of missing exact
certificates for the sixth-power theorem. The earlier integrated
collection remains the record of the other complete proofs,
candidates, and counterexamples.

\clearpage\appendix
\section{Standalone exact sixth-power certificate}
\label{app:certificate}
This complete Python program proves the finite statements used in
Section~\ref{war:section}. In the embedded data, a pair $a:c$ means
that the positive sixth power $a^6$ occurs with multiplicity $c$.
All seventy-three rows are present. Integer bit operations encode
the full finite sets; no floating-point arithmetic is used.

\begin{Verbatim}[fontsize=\footnotesize,breaklines,breakanywhere]
#!/usr/bin/env python3
"""Replay all finite certificates for exactly 78 positive sixth powers.

Python 3.10+; standard Python only; no input files or downloads.
Run normally (without -O) so every assertion is active.
The infinite proof additionally uses the classical theorem g(6)=73.
A row s lists base:multiplicity pairs for a representation of P.
"""

P = 41_253_168_892
H = 3_500_000
WITNESSES = {
    5: "10:1 13:1 45:1 47:1 53:1",
    6: "11:1 23:1 25:1 32:1 42:1 57:1",
    7: "7:1 13:1 24:1 25:1 36:1 49:1 54:1",
    8: "21:2 36:1 37:1 38:1 44:1 48:1 49:1",
    9: "7:1 9:1 11:1 24:1 28:1 30:1 39:1 42:1 56:1",
    10: "3:1 6:1 7:1 11:1 21:1 28:1 36:2 42:1 56:1",
    11: "3:1 6:1 14:3 21:2 39:1 42:1 44:1 54:1",
    12: "5:1 12:3 18:1 19:1 23:1 24:2 27:1 36:1 58:1",
    13: "2:1 8:1 14:1 17:2 19:1 22:3 26:1 31:1 34:1 58:1",
    14: "2:2 4:2 5:1 7:1 8:2 9:1 14:1 16:1 23:1 38:1 58:1",
    15: "3:2 5:1 7:3 10:1 13:5 23:1 38:1 58:1",
    16: "2:1 3:1 6:2 7:1 8:1 11:2 14:3 16:2 22:1 38:1 58:1",
    17: "2:1 6:2 7:2 8:1 10:1 12:1 13:1 14:3 15:1 16:1 22:1 38:1 58:1",
    18: "1:3 2:2 3:2 4:1 7:4 9:1 10:1 21:2 38:1 58:1",
    19: "3:6 6:2 9:5 12:2 16:1 23:1 38:1 58:1",
    20: "1:2 3:3 4:1 5:3 6:3 7:1 8:2 9:1 17:1 23:1 38:1 58:1",
    21: "1:2 3:7 4:3 5:1 6:1 7:1 10:1 14:1 16:1 23:1 38:1 58:1",
    22: "1:1 3:6 5:2 6:3 7:2 8:3 14:1 16:1 23:1 38:1 58:1",
    23: "1:2 2:5 5:4 6:1 7:4 8:3 17:1 23:1 38:1 58:1",
    24: "2:3 3:6 4:1 6:4 7:4 9:1 14:1 16:1 23:1 38:1 58:1",
    25: "1:1 2:8 4:1 5:2 6:2 7:6 9:1 17:1 23:1 38:1 58:1",
    26: "1:4 2:7 3:3 4:2 5:3 6:1 8:1 10:1 17:1 23:1 38:1 58:1",
    27: "2:1 3:1 4:15 6:4 7:1 10:1 17:1 23:1 38:1 58:1",
    28: "1:2 2:5 3:5 4:4 5:1 6:1 7:2 8:4 17:1 23:1 38:1 58:1",
    29: "1:8 2:4 3:7 4:2 7:3 10:1 17:1 23:1 38:1 58:1",
    30: "1:4 2:5 3:6 4:1 5:3 6:1 7:4 8:1 9:1 17:1 23:1 38:1 58:1",
    31: "2:4 3:6 5:7 6:3 7:5 8:2 17:1 23:1 38:1 58:1",
    32: "1:11 2:1 4:1 5:8 7:6 9:1 17:1 23:1 38:1 58:1",
    33: "1:9 2:6 3:9 4:1 6:2 8:1 10:1 17:1 23:1 38:1 58:1",
    34: "1:2 2:6 3:9 4:2 5:6 6:3 7:1 10:1 17:1 23:1 38:1 58:1",
    35: "1:4 2:17 3:2 4:1 5:2 6:2 7:2 10:1 17:1 23:1 38:1 58:1",
    36: "1:6 2:9 4:8 6:4 7:3 8:1 9:1 17:1 23:1 38:1 58:1",
    37: "1:9 2:4 3:12 6:2 7:4 8:1 9:1 17:1 23:1 38:1 58:1",
    38: "1:11 2:8 3:1 4:3 6:5 7:5 9:1 17:1 23:1 38:1 58:1",
    39: "1:4 2:15 3:5 5:2 6:2 7:6 9:1 17:1 23:1 38:1 58:1",
    40: "1:7 2:14 3:8 4:1 5:3 6:1 8:1 10:1 17:1 23:1 38:1 58:1",
    41: "1:7 2:5 3:15 4:1 5:3 6:4 7:1 10:1 17:1 23:1 38:1 58:1",
    42: "1:21 2:1 3:9 5:2 6:2 7:2 10:1 17:1 23:1 38:1 58:1",
    43: "1:11 2:11 3:12 4:1 7:3 10:1 17:1 23:1 38:1 58:1",
    44: "1:28 2:1 3:1 4:2 6:2 7:4 8:1 9:1 17:1 23:1 38:1 58:1",
    45: "1:9 2:16 4:3 5:3 6:4 7:5 9:1 17:1 23:1 38:1 58:1",
    46: "1:8 2:14 3:3 4:8 6:2 7:6 9:1 17:1 23:1 38:1 58:1",
    47: "1:12 2:13 3:14 6:2 8:1 10:1 17:1 23:1 38:1 58:1",
    48: "1:14 2:17 3:3 4:3 6:5 7:1 10:1 17:1 23:1 38:1 58:1",
    49: "1:7 2:24 3:7 5:2 6:2 7:2 10:1 17:1 23:1 38:1 58:1",
    50: "1:30 2:8 3:1 4:3 7:3 10:1 17:1 23:1 38:1 58:1",
    51: "1:26 2:9 4:2 5:3 6:1 7:4 8:1 9:1 17:1 23:1 38:1 58:1",
    52: "1:14 2:15 3:6 4:2 6:5 7:5 9:1 17:1 23:1 38:1 58:1",
    53: "1:16 2:10 3:10 4:3 5:1 6:2 7:6 9:1 17:1 23:1 38:1 58:1",
    54: "1:31 2:10 3:3 4:2 6:2 8:1 10:1 17:1 23:1 38:1 58:1",
    55: "2:40 3:1 4:3 6:5 7:1 10:1 17:1 23:1 38:1 58:1",
    56: "1:2 2:35 3:5 4:4 5:1 6:2 7:2 10:1 17:1 23:1 38:1 58:1",
    57: "1:14 2:18 3:17 7:3 10:1 17:1 23:1 38:1 58:1",
    58: "1:31 2:8 3:6 4:1 6:2 7:4 8:1 9:1 17:1 23:1 38:1 58:1",
    59: "2:38 3:4 4:2 6:5 7:5 9:1 17:1 23:1 38:1 58:1",
    60: "1:2 2:33 3:8 4:3 5:1 6:2 7:6 9:1 17:1 23:1 38:1 58:1",
    61: "1:17 2:33 3:1 4:2 6:2 8:1 10:1 17:1 23:1 38:1 58:1",
    62: "1:17 2:24 3:8 4:2 6:5 7:1 10:1 17:1 23:1 38:1 58:1",
    63: "1:43 2:5 3:3 4:1 5:2 6:2 7:2 10:1 17:1 23:1 38:1 58:1",
    64: "1:33 2:15 3:6 4:2 7:3 10:1 17:1 23:1 38:1 58:1",
    65: "1:17 2:31 3:4 4:1 6:2 7:4 8:1 9:1 17:1 23:1 38:1 58:1",
    66: "1:17 2:22 3:11 4:1 6:5 7:5 9:1 17:1 23:1 38:1 58:1",
    67: "1:43 2:3 3:6 5:2 6:2 7:6 9:1 17:1 23:1 38:1 58:1",
    68: "1:34 2:17 3:8 4:1 6:2 8:1 10:1 17:1 23:1 38:1 58:1",
    69: "1:3 2:47 3:6 4:2 6:5 7:1 10:1 17:1 23:1 38:1 58:1",
    70: "1:29 2:28 3:1 4:1 5:2 6:2 7:2 10:1 17:1 23:1 38:1 58:1",
    71: "1:19 2:38 3:4 4:2 7:3 10:1 17:1 23:1 38:1 58:1",
    72: "1:3 2:54 3:2 4:1 6:2 7:4 8:1 9:1 17:1 23:1 38:1 58:1",
    73: "1:36 2:19 4:3 6:5 7:5 9:1 17:1 23:1 38:1 58:1",
    74: "1:29 2:26 3:4 5:2 6:2 7:6 9:1 17:1 23:1 38:1 58:1",
    75: "1:20 2:40 3:6 4:1 6:2 8:1 10:1 17:1 23:1 38:1 58:1",
    76: "1:53 2:5 3:4 4:3 6:5 7:1 10:1 17:1 23:1 38:1 58:1",
    77: "1:46 2:12 3:8 5:2 6:2 7:2 10:1 17:1 23:1 38:1 58:1",
}

# Elementary Apery representatives of <63,728,117648>, modulo 63.
assert 9 * 728 == 104 * 63
assert 7 * 117648 == 13072 * 63
least = [None] * 63
for y in range(9):
    for z in range(7):
        d = 728 * y + 117648 * z
        r = d % 63
        assert least[r] is None
        least[r] = d
assert all(d is not None for d in least)
assert least[0] == 0
largest_gap = max(least) - 63
number_of_gaps = sum((d - r) // 63 for r, d in enumerate(least))
assert largest_gap == 711649
assert number_of_gaps == 355825
# All generators a^6-1 with a>=10 exceed the Frobenius number.
assert 10**6 - 1 > largest_gap
# Every remaining sixth-power difference is already in this semigroup.
assert all(a**6 - 1 >= least[(a**6 - 1) % 63]
           for a in range(1, 10))
print("Permanent gaps:", number_of_gaps, "largest:", largest_gap)

# Truncated nonnegative sumsets: exact membership, not an approximation.
mask = (1 << (H + 1)) - 1
coins = [a**6 - 1 for a in range(1, 13)]
assert coins[-1] <= H < 13**6 - 1
v = 1
recorded = {}
for j in range(1, 79):
    w = 0
    for c in coins:
        w |= v << c
    v = w & mask
    if j in (40, 77, 78):
        gaps = mask ^ v
        recorded[j] = (gaps.bit_count(), gaps.bit_length() - 1)
        print("Row", j, "missing count and largest:", recorded[j])
    if j == 77:
        assert not ((v >> 715744) & 1)
assert recorded == {
    40: (360244, 746110),
    77: (355826, 715744),
    78: (355825, 711649),
}

# Independently form all semigroup membership bits and compare every bit.
semigroup = bytearray((H + 8) // 8)
for d in least:
    for n in range(d, H + 1, 63):
        semigroup[n >> 3] |= 1 << (n & 7)
assert int.from_bytes(semigroup, "little") == v
assert 715744 >= least[715744 % 63]
print("Every membership bit at row 78 agrees with the semigroup.")
print("715744 belongs to the semigroup but does not belong to 77D.")

# Six exact interval extensions, starting from row 40.
lo, hi = 746111, H
expected_ends = (
    11029535, 45041758, 236144733,
    1527612701, 13758203164, 177964693339,
)
for j, a, expected in zip(
        range(41, 47), (14, 18, 24, 33, 48, 74), expected_ends):
    assert a**6 - (a - 1)**6 <= hi - lo + 1
    hi += a**6 - 1
    assert hi == expected
    print("Interval", j, "largest appended base", a, ":", lo, hi)
assert lo <= H and hi >= P - 78

# Every embedded padding certificate is checked by exact integer sums.
assert set(WITNESSES) == set(range(5, 78))
for s, encoded in WITNESSES.items():
    pairs = [tuple(map(int, item.split(":")))
             for item in encoded.split()]
    assert all(len(pair) == 2 for pair in pairs)
    assert all(a > 0 and count > 0 for a, count in pairs)
    assert len({a for a, count in pairs}) == len(pairs)
    assert sum(count for a, count in pairs) == s
    assert sum(count * a**6 for a, count in pairs) == P
print("All 73 embedded padding representations passed exact verification.")
print("PASS: all finite certificates for the global proof are verified.")
\end{Verbatim}

\clearpage
\section{Archive locations and reproducibility record}
\label{app:archive}
The source checkpoints are the following files in the supplied
archive, relative to its expanded root. Their earlier
limitations are stated in their own prose.
\begin{itemize}\small\raggedright
\item Round 12: \path{Tournament_Round_12/heisei/manuscript.tex}.
\item Round 22: \path{Tournament_Round_22_Final/Tournament_Round_22/round22/taisho/taisho.tex}.
\item Round 25: \path{Tournament_Round_25_Final/Tournament_Round_25/round25/taisho/paper.tex}.
\item Round 21: \path{Tournament_Round_21_Final/Tournament_Round_21/round21/keio/paper.tex}.
\end{itemize}
The exact verifier reports the finite rows
$(360244,746110)$, $(355826,715744)$ and $(355825,711649)$
at indices $40,77,78$, respectively; compares every membership
bit at index $78$; verifies all six interval extensions; and
verifies all seventy-three padding identities. These are the
finite premises used in the proof.

The analytic scale lemma uses the rational margins
$3161383/510000000$ and $47/82620$ above $13/20$.
The sine-quotient proof uses the uniform derivative estimate
$J_q'\le1/25$ only on $u\le1$, where its majorant converges.
No empirical cutoff is used to infer an infinite analytic claim.

\Needspace{10\baselineskip}
\phantomsection
\addcontentsline{toc}{section}{References}
\begin{thebibliography}{9}
\bibitem{Cazenave2026}
T. Cazenave,
\emph{Lemmas and Conjectures in Spectral Graph Theory},
arXiv:2611.03879v1, 2 November 2026.
\url{https://arxiv.org/abs/2611.03879v1}.
Section 7, Conjecture 6(iii).

\bibitem{AlzerKwong}
H. Alzer and M. K. Kwong,
\emph{On a Sine Series of Ramanujan},
Axioms \textbf{15} (2026), no. 11, article 859.
\href{https://doi.org/10.3390/axioms15110859}{doi:10.3390/axioms15110859}.
Section 4, conjecture (27).

\bibitem{Malesevic}
B. Male\v sevi\'c, M. Mi\'covi\'c, B. Mihailovi\'c and T. Lutovac,
\emph{A new approach to selecting constants for some analytic inequalities},
Journal of Mathematical Inequalities \textbf{19} (2025), 223--245.
\href{https://doi.org/10.7153/jmi-2025-19-15}{doi:10.7153/jmi-2025-19-15}.
Equation (15) and Section 5.

\bibitem{HungLi}
L.-C. Hung and P.-Y. Li,
\emph{On generalization of D'Aurizio--S\'andor inequalities involving a parameter},
Journal of Mathematical Inequalities \textbf{12} (2018), 853--860.
\href{https://doi.org/10.7153/jmi-2018-12-63}{doi:10.7153/jmi-2018-12-63}.

\bibitem{BenfieldLippard}
B. Benfield and O. Lippard,
\emph{Integers that are not the Sum of Positive Powers},
arXiv:2404.08193v2, 31 March 2025.
\url{https://arxiv.org/abs/2404.08193v2}.
Conjecture 6.4, Table 1, and Section 11.

\bibitem{Pillai1940}
S. S. Pillai,
\emph{On Waring's problem $g(6)=73$},
Proceedings of the Indian Academy of Sciences, Section A
\textbf{12} (1940), 30--40.
\href{https://doi.org/10.1007/BF03170721}{doi:10.1007/BF03170721}.
\end{thebibliography}
\end{document}
